\section{上极限与下极限}
\label{sec:1.5}

这一节介绍上极限与下极限的概念，并应用其是否相等来刻画数列有没有极限.

\begin{definition}\label{def:1.5.1}
如果 $a \in \mathbb{R}_* = \mathbb{R} \cup \{+\infty, -\infty\}$ 的任何开邻域内都含数列 $\{x_n\}$ 的无限项，则称 $a$ 为数列 $\{x_n\}$ 的一个\textbf{聚点}。易知，数列 $\{x_n\}$ 的聚点实际上就是 $\{x_n\}$ 的某个子列的极限。
\end{definition}

显然，数集 $\{x_n \mid n \in \mathbb{N}\}$ 的聚点 $a$ 必为数列 $\{x_n\}$ 的聚点。但反之不真，见例~\ref{ex:1.5.1}(1)(4)。

\begin{definition}\label{def:1.5.2}
数列 $\{x_n\}$ 的最大聚点 $a_{\text{大}}$ 与最小聚点 $a_{\text{小}}$（由定理~\ref{thm:1.5.1} 知，它们一定存在）分别称为数列 $\{x_n\}$ 的\textbf{上极限}与\textbf{下极限}，记作
\[
a_{\text{大}} = \varlimsup_{n\to+\infty} x_n,\quad a_{\text{小}} = \varliminf_{n\to+\infty} x_n.
\]
由定义，显然有
\[
a_{\text{小}} = \varliminf_{n\to+\infty} x_n \leqslant \varlimsup_{n\to+\infty} x_n = a_{\text{大}}.
\]
\end{definition}

\begin{example}\label{ex:1.5.1}
\begin{enumerate}
    \item 常数列 $\{a, a, a, \cdots\}$ 只有一个聚点 $a$，但其对应的数集 $\{a\}$ 无聚点。
    \item 数列 $\left\{\frac{1}{n}\right\}$ 只有一个聚点 $0$，其对应的数集 $\left\{\frac{1}{n} \;\middle|\; n \in \mathbb{N}\right\}$ 的聚点也只有一个 $0$。
    \item 数列 $\left\{(-1)^n \frac{n}{n+1}\right\}$ 的聚点集为 $\{-1, 1\}$，其对应的数集 $\left\{(-1)^n \frac{n}{n+1} \;\middle|\; n \in \mathbb{N}\right\}$ 的聚点集也为 $\{-1, 1\}$。
    \item 数列 $\left\{\sin\frac{n\pi}{4}\right\}$ 的聚点集为 $\left\{-1, -\frac{\sqrt{2}}{2}, 0, \frac{\sqrt{2}}{2}, 1\right\}$，其对应的数集 $\left\{\sin\frac{n\pi}{4} \;\middle|\; n \in \mathbb{N}\right\} = \left\{-1, -\frac{\sqrt{2}}{2}, 0, \frac{\sqrt{2}}{2}, 1\right\}$ 无聚点。
\end{enumerate}
易见，
\[
\varlimsup_{n\to+\infty} a = a = \varliminf_{n\to+\infty} a;
\]
\[
\varlimsup_{n\to+\infty} \frac{1}{n} = 0 = \varliminf_{n\to+\infty} \frac{1}{n};
\]
\[
\varlimsup_{n\to+\infty} (-1)^n \frac{n}{n+1} = 1,\quad \varliminf_{n\to+\infty} (-1)^n \frac{n}{n+1} = -1;
\]
\[
\varlimsup_{n\to+\infty} \sin\frac{n\pi}{4} = 1,\quad \varliminf_{n\to+\infty} \sin\frac{n\pi}{4} = -1.
\]
\end{example}

\begin{example}\label{ex:1.5.2}
\begin{enumerate}
    \item $\varlimsup\limits_{n\to+\infty} \left[(-1)^n + 1\right]n = +\infty,\quad \varliminf\limits_{n\to+\infty} \left[(-1)^n + 1\right]n = 0$;
    \item $\varlimsup\limits_{n\to+\infty} n = +\infty = \varliminf\limits_{n\to+\infty} n$;
    \item $\varlimsup\limits_{n\to+\infty} \left[(-1)^{n+1} - 1\right]n = 0,\quad \varliminf\limits_{n\to+\infty} \left[(-1)^{n+1} - 1\right]n = -\infty$;
    \item $\varlimsup\limits_{n\to+\infty} (-n^2) = -\infty = \varliminf\limits_{n\to+\infty} (-n^2)$;
    \item $\varlimsup\limits_{n\to+\infty} (-1)^{n-1} = 1,\quad \varliminf\limits_{n\to+\infty} (-1)^{n-1} = -1$.
\end{enumerate}
\end{example}

\begin{theorem}\label{thm:1.5.1}
\begin{enumerate}
    \item 有界数列 $\{x_n\}$ 至少有一个聚点，存在最大聚点与最小聚点，且这两个聚点都为实数，它们分别为上极限 $\varlimsup\limits_{n\to+\infty} x_n$ 与下极限 $\varliminf\limits_{n\to+\infty} x_n$。
    \item 如果数列 $\{x_n\}$ 无上界，则 $\varlimsup\limits_{n\to+\infty} x_n = +\infty$，此时 $+\infty$ 为数列 $\{x_n\}$ 的最大聚点。
    
    如果数列 $\{x_n\}$ 有上界 $b$：
    \begin{enumerate}
        \item[①] 若 $\forall a < b$, $[a, b]$ 中只含数列 $\{x_n\}$ 的有限项，则 $\lim\limits_{n\to+\infty} x_n = -\infty = \varlimsup\limits_{n\to+\infty} x_n$（此时，$\varliminf\limits_{n\to+\infty} x_n = -\infty$）。
        \item[②] 若 $\exists a < b$, $[a, b]$ 中含数列 $\{x_n\}$ 的无限项，则数列 $\{x_n\}$ 以实数为最大聚点，它就是 $\varlimsup\limits_{n\to+\infty} x_n$。
    \end{enumerate}
    \item 如果数列 $\{x_n\}$ 无下界，则 $\varliminf\limits_{n\to+\infty} x_n = -\infty$，此时 $-\infty$ 为数列 $\{x_n\}$ 的最小聚点。
    
    如果数列 $\{x_n\}$ 有下界 $b$：
    \begin{enumerate}
        \item[①] 若 $\forall b > a$, $[a, b]$ 中只含数列 $\{x_n\}$ 的有限项，则 $\lim\limits_{n\to+\infty} x_n = +\infty = \varliminf\limits_{n\to+\infty} x_n$（此时，$\varlimsup\limits_{n\to+\infty} x_n = +\infty$）。
        \item[②] 若 $\exists b > a$, $[a, b]$ 中含数列 $\{x_n\}$ 的无限项，则数列 $\{x_n\}$ 以实数为最小聚点，它就是 $\varliminf\limits_{n\to+\infty} x_n$。
    \end{enumerate}
\end{enumerate}
\end{theorem}

\begin{proof}
(1) 因数列 $\{x_n\}$ 有界，令 $\{x_n \mid n \in \mathbb{N}\} \subset [-M, M] = [a_1, b_1]$。将 $[a_1, b_1]$ 两等分，则必有一等分含数列 $\{x_n\}$ 的无限多项，记此子区间为 $[a_2, b_2]$，则 $[a_1, b_1] \supset [a_2, b_2]$，且
\[
b_2 - a_2 = \frac{1}{2}(b_1 - a_1) = M.
\]
再将 $[a_2, b_2]$ 两等分，则必有一等分含数列 $\{x_n\}$ 的无穷多项，记此子区间为 $[a_3, b_3]$，则 $[a_2, b_2] \supset [a_3, b_3]$，且
\[
b_3 - a_3 = \frac{1}{2}(b_2 - a_2) = \frac{M}{2}.
\]
如此下去，得到一个递降闭区间套：
\[
[a_1, b_1] \supset [a_2, b_2] \supset \cdots \supset [a_k, b_k] \supset \cdots,
\]
\[
b_k - a_k = \frac{M}{2^{k-1}} \to 0\quad (k \to +\infty),
\]
且每个闭区间 $[a_k, b_k]$ 都含数列 $\{x_n\}$ 的无限多项。

由实数连续性命题（三）（闭区间套原理）知，$\exists_1 x_0 \in \bigcap_{k=1}^\infty [a_k, b_k]$。对 $x_0$ 的任何开邻域 $U$，$\exists \varepsilon > 0$, s.t. $B(x_0; \varepsilon) = (x_0 - \varepsilon, x_0 + \varepsilon) \subset U$。则 $\exists N \in \mathbb{N}$，当 $k > N$ 时，$[a_k, b_k] \subset (x_0 - \varepsilon, x_0 + \varepsilon) \subset U$，从而 $U$ 中含数列 $\{x_n\}$ 的无限多项，所以 $x_0$ 为数列 $\{x_n\}$ 的聚点。

至于最大聚点的存在性，只需在上述证明过程中，当每次将区间 $[a_{k-1}, b_{k-1}]$ 等分为两个子区间时，若右边一个含数列 $\{x_n\}$ 的无限项，则取它为 $[a_k, b_k]$；若右边一个含数列 $\{x_n\}$ 的有限项，则取左边的子区间为 $[a_k, b_k]$。于是，所选 $[a_k, b_k]$ 都含数列 $\{x_n\}$ 的无限项，同时在 $[a_k, b_k]$ 的右边却至多含数列 $\{x_n\}$ 的有限项，其中
\[
b_k - a_k = \frac{1}{2}(b_{k-1} - a_{k-1}) = \cdots = \frac{1}{2^{k-1}}(b_1 - a_1) \to 0\quad (k \to +\infty).
\]
再根据闭区间套原理知，$\exists_1 x_0 \in \bigcap_{k=1}^\infty [a_k, b_k]$。下证 $x_0$ 为数列 $\{x_n\}$ 的最大聚点。（反证）若不然，设另有数列 $\{x_n\}$ 的聚点 $x_0^* > x_0$。令 $\delta = \frac{1}{3}(x_0^* - x_0) > 0$，则在 $B(x_0^*; \delta) = (x_0^* - \delta, x_0^* + \delta)$ 内含数列 $\{x_n\}$ 的无限多项。但当 $k$ 充分大时，$B(x_0^*; \delta)$ 完全落在 $[a_k, b_k]$ 的右边。这与上述 $[a_k, b_k]$ 右边至多含数列 $\{x_n\}$ 的有限项相矛盾。

类似可证最小聚点的存在性，或用 $\{-x_n\}$ 代替 $\{x_n\}$。

(2) 如果数列 $\{x_n\}$ 无上界，根据例~\ref{ex:1.1.16} 可知，$\{x_n\}$ 有子列 $\{x_{n_k}\}$，s.t. $\lim\limits_{k\to+\infty} x_{n_k} = +\infty$。因此，$+\infty$ 为数列 $\{x_n\}$ 最大聚点，从而 $\varlimsup\limits_{n\to+\infty} x_n = +\infty$。

如果数列 $\{x_n\}$ 有上界 $b$：
\begin{enumerate}
    \item[①] 若 $\forall a < b$, $[a, b]$ 中只含数列 $\{x_n\}$ 的有限项。则根据极限为 $-\infty$ 的定义立知，
    \[
    \lim_{n\to+\infty} x_n = -\infty = \varlimsup_{n\to+\infty} x_n.
\]
    \item[②] 若 $\exists a < b$, $[a, b]$ 中含数列 $\{x_n\}$ 的无限项。由 (1) 的结果，数列 $\{x_n\} \cap [a, b]$ 有最大聚点，显然它也是数列 $\{x_n\}$ 的最大聚点，即为 $\varlimsup\limits_{n\to+\infty} x_n$。
\end{enumerate}

(3) 类似 (2) 证明，或用 $\{-x_n\}$ 代替 $\{x_n\}$。
\end{proof}

\begin{theorem}\label{thm:1.5.2}
\[
\lim_{n\to+\infty} x_n = a \iff \varlimsup_{n\to+\infty} x_n = \varliminf_{n\to+\infty} x_n = a.
\]
\end{theorem}

\begin{proof}
($\Rightarrow$) 设 $\lim\limits_{n\to+\infty} x_n = a$，则对 $a$ 的任一开邻域 $U$，$\exists N \in \mathbb{N}$，当 $n > N$ 时，$x_n \in U$，从而 $a$ 为数列 $\{x_n\}$ 的一个聚点。

$\forall b \neq a$，则存在 $a$ 的开邻域 $U_a$ 与 $b$ 的开邻域 $U_b$，s.t. $U_a \cap U_b = \varnothing$。由于 $\lim\limits_{n\to+\infty} x_n = a$，故 $\exists N \in \mathbb{N}$，当 $n > N$ 时，$x_n \in U_a$，所以 $x_n \notin U_b$，从而 $U_b$ 中至多含数列 $\{x_n\}$ 的有限项（如 $x_1, \cdots, x_N$）。因此，$b$ 不为数列 $\{x_n\}$ 的聚点。

综上可知，$a$ 为数列 $\{x_n\}$ 的惟一聚点。所以，
\[
\varlimsup_{n\to+\infty} x_n = a = \varliminf_{n\to+\infty} x_n.
\]
或者，因 $\lim\limits_{n\to+\infty} x_n = a$，故 $\{x_n\}$ 的任何子列 $\{x_{n_k}\}$ 也必有 $\lim\limits_{k\to+\infty} x_{n_k} = a$。因此，数列 $\{x_n\}$ 有唯一的聚点 $a$，从而
\[
\varlimsup_{n\to+\infty} x_n = a = \varliminf_{n\to+\infty} x_n.
\]

($\Leftarrow$) 设 $\varlimsup\limits_{n\to+\infty} x_n = \varliminf\limits_{n\to+\infty} x_n = a$，则数列 $\{x_n\}$ 只有一个聚点 $a$。因此，对 $a$ 的任何开邻域 $U$，在 $U$ 外只含数列 $\{x_n\}$ 的有限项 $x_{n_1}, \cdots, x_{n_k}$（否则数列 $\{x_n\}$ 在 $U$ 外还有异于 $a$ 的聚点，这与数列 $\{x_n\}$ 只有一个聚点相矛盾）。于是，当 $n > N = \max\{n_1, \cdots, n_k, 1\}$ 时，有 $x_n \in U$。这就证明了 $\lim\limits_{n\to+\infty} x_n = a$。
\end{proof}

\begin{theorem}\label{thm:1.5.3}
设 $\{x_n\}$ 为有界数列，则下列结论等价：
\begin{enumerate}
    \item $a_{\text{大}}$ 为数列 $\{x_n\}$ 的上极限。
    \item $\forall \varepsilon > 0$, $\exists N \in \mathbb{N}$, s.t. 当 $n > N$ 时，有 $x_n < a_{\text{大}} + \varepsilon$；且 $\exists$ 子列 $\{x_{n_k}\}$, s.t. $x_{n_k} > a_{\text{大}} - \varepsilon$, $\forall k \in \mathbb{N}$。
    \item $\forall \alpha > a_{\text{大}}$，数列 $\{x_n\}$ 中大于 $\alpha$ 的项至多有限个；$\forall \beta < a_{\text{大}}$，数列 $\{x_n\}$ 中大于 $\beta$ 的项有无限多个。
\end{enumerate}
\end{theorem}

\begin{proof}
(1) $\Rightarrow$ (2)

因 $a_{\text{大}}$ 为数列 $\{x_n\}$ 的聚点，故 $\forall \varepsilon > 0$，在 $B(a_{\text{大}}; \varepsilon) = (a_{\text{大}} - \varepsilon, a_{\text{大}} + \varepsilon)$ 内含数列 $\{x_n\}$ 的无限多项 $\{x_{n_k} \mid n_1 < n_2 < \cdots\}$，则有 $x_{n_k} > a_{\text{大}} - \varepsilon$, $\forall k \in \mathbb{N}$。

又因 $a_{\text{大}}$ 为数列 $\{x_n\}$ 的最大聚点，故在 $a_{\text{大}} + \varepsilon$ 的右边至多只有数列 $\{x_n\}$ 的有限项（否则必有数列 $\{x_n\}$ 的聚点 $\geqslant a_{\text{大}} + \varepsilon$，这与 $a_{\text{大}}$ 为数列 $\{x_n\}$ 的最大聚点相矛盾）。设此有限项的最大指标为 $N$，则当 $n > N$ 时，有 $x_n < a_{\text{大}} + \varepsilon$。

(2) $\Rightarrow$ (3)

$\forall \alpha > a_{\text{大}}$，令 $\varepsilon = \alpha - a_{\text{大}} > 0$。由 (2) 知，$\exists N \in \mathbb{N}$，当 $n > N$ 时，有 $x_n < a_{\text{大}} + \varepsilon = a_{\text{大}} + (\alpha - a_{\text{大}}) = \alpha$。故数列 $\{x_n\}$ 中大于 $\alpha$ 的项至多有限个。

$\forall \beta < a_{\text{大}}$，令 $\varepsilon = a_{\text{大}} - \beta > 0$。由 (2) 知，存在 $\{x_n\}$ 的子列 $\{x_{n_k}\}$，s.t. $x_{n_k} > a_{\text{大}} - \varepsilon = \beta$, $\forall k \in \mathbb{N}$。故数列 $\{x_n\}$ 中大于 $\beta$ 的项有无限多个。

(3) $\Rightarrow$ (1)

设 $U$ 为 $a_{\text{大}}$ 的任一开邻域，则 $\exists \varepsilon > 0$, s.t. $B(a_{\text{大}}; \varepsilon) = (a_{\text{大}} - \varepsilon, a_{\text{大}} + \varepsilon) \subset U$。由于 $\alpha = a_{\text{大}} + \varepsilon > a_{\text{大}}$，根据 (3)，$\{x_n\}$ 中大于 $\alpha = a_{\text{大}} + \varepsilon$ 至多有限个。而对 $\beta = a_{\text{大}} - \varepsilon < a_{\text{大}}$，根据 (3)，$\{x_n\}$ 中大于 $\beta = a_{\text{大}} - \varepsilon$ 有无限项。因此，$(a_{\text{大}} - \varepsilon, a_{\text{大}} + \varepsilon)$ 中含 $\{x_n\}$ 的无限项，从而 $U$ 中含 $\{x_n\}$ 的无限项。这就证明了 $a_{\text{大}}$ 为数列 $\{x_n\}$ 的一个聚点。

另一方面，$\forall \alpha > a_{\text{大}}$，记 $\varepsilon = \frac{1}{2}(\alpha - a_{\text{大}})$。由 (3) 可知，$\{x_n\}$ 中大于 $a_{\text{大}} + \varepsilon\;(> a_{\text{大}})$ 的项至多有限个。故 $\alpha$ 不为数列 $\{x_n\}$ 的聚点。这就证明了 $a_{\text{大}}$ 为数列 $\{x_n\}$ 的最大聚点，即 $a_{\text{大}}$ 为数列 $\{x_n\}$ 的上极限。
\end{proof}

\begin{theorem}\label{thm:1.5.4}
设 $\{x_n\}$ 为有界数列，则下列结论等价：
\begin{enumerate}
    \item $a_{\text{小}}$ 为数列 $\{x_n\}$ 的下极限。
    \item $\forall \varepsilon > 0$, $\exists N \in \mathbb{N}$, s.t. 当 $n > N$ 时，有 $x_n > a_{\text{小}} - \varepsilon$；且 $\exists$ 子列 $\{x_{n_k}\}$, s.t. $x_{n_k} < a_{\text{小}} + \varepsilon$, $\forall k \in \mathbb{N}$。
    \item $\forall \beta < a_{\text{小}}$，数列 $\{x_n\}$ 中小于 $\beta$ 的项至多有限个；$\forall \alpha > a_{\text{小}}$，数列 $\{x_n\}$ 中小于 $\alpha$ 的项有无限多个。
\end{enumerate}
\end{theorem}

\begin{proof}
类似定理~\ref{thm:1.5.3} 证明，或用 $\{-x_n\}$ 代替 $\{x_n\}$。
\end{proof}

\begin{theorem}\label{thm:1.5.5}
\begin{enumerate}
    \item \textbf{（上、下极限的保不等式性）} 设数列 $\{a_n\}, \{b_n\}$ 满足：$\exists N_0 \in \mathbb{N}$，当 $n > N_0$ 时，有 $a_n \leqslant b_n$，则
    \[
    \varlimsup_{n\to+\infty} a_n \leqslant \varlimsup_{n\to+\infty} b_n,\quad \varliminf_{n\to+\infty} a_n \leqslant \varliminf_{n\to+\infty} b_n.
\]
    特别地，若 $\alpha, \beta$ 为常数，又 $\exists N_0 \in \mathbb{N}$，当 $n > N_0$ 时，有 $\alpha \leqslant a_n \leqslant \beta$，则
    \[
    \alpha \leqslant \varliminf_{n\to+\infty} a_n \leqslant \varlimsup_{n\to+\infty} a_n \leqslant \beta.
\]
    \item 设 $\{a_n\}, \{b_n\}$ 为有界数列，则
    \[
    \varlimsup_{n\to+\infty} (a_n + b_n) \leqslant \varlimsup_{n\to+\infty} a_n + \varlimsup_{n\to+\infty} b_n.
\]
    举例说明严格不等号可能成立：
    \[
    \varlimsup_{n\to+\infty} (a_n + b_n) < \varlimsup_{n\to+\infty} a_n + \varlimsup_{n\to+\infty} b_n.
\]
\end{enumerate}
\end{theorem}

\begin{proof}
(1) 如果 $a_{\text{大}} = \varlimsup\limits_{n\to+\infty} a_n = +\infty$，则数列 $\{a_n\}$ 无上界。由 $a_n \leqslant b_n\;(n > N_0)$，故 $\{b_n\}$ 也无上界，从而
\[
a_{\text{大}} = \varlimsup_{n\to+\infty} a_n = +\infty = \varlimsup_{n\to+\infty} b_n = b_{\text{大}}.
\]
如果 $a_{\text{大}} = \varlimsup\limits_{n\to+\infty} a_n = -\infty$，则显然 $\varlimsup\limits_{n\to+\infty} a_n = -\infty \leqslant \varlimsup\limits_{n\to+\infty} b_n$。

如果 $-\infty < a_{\text{大}} = \varlimsup\limits_{n\to+\infty} a_n < +\infty$。根据定理~\ref{thm:1.5.3}(2)，$\exists \{a_n\}$ 的子列 $\{a_{n_k}\}$，s.t. $a_{n_k} > a_{\text{大}} - \varepsilon$, $\forall k \in \mathbb{N}$。故由题设可知，
\[
b_{n_k} \geqslant a_{n_k} > a_{\text{大}} - \varepsilon,\quad n_k > N_0.
\]
于是，数列 $\{b_n\}$ 必有聚点大于 $a_{\text{大}} - \varepsilon$，由此知 $\{b_n\}$ 的最大聚点（上极限）
\[
b_{\text{大}} = \varlimsup_{n\to+\infty} b_n \geqslant a_{\text{大}} - \varepsilon.
\]
令 $\varepsilon \to 0^+$，得到 $\varlimsup\limits_{n\to+\infty} b_n \geqslant a_{\text{大}} = \varlimsup\limits_{n\to+\infty} a_n$。

类似可证 $\varliminf\limits_{n\to+\infty} a_n \leqslant \varliminf\limits_{n\to+\infty} b_n$。

或者，因 $a_n \leqslant b_n$，故 $-a_n \geqslant -b_n$，由上述结论得到
\[
-\varliminf_{n\to+\infty} b_n = \varlimsup_{n\to+\infty} (-b_n) \leqslant \varlimsup_{n\to+\infty} (-a_n) = -\varliminf_{n\to+\infty} a_n,
\]
\[
\varliminf_{n\to+\infty} b_n \geqslant \varliminf_{n\to+\infty} a_n.
\]
当 $\alpha \leqslant a_n \leqslant \beta$ 时，
\[
\alpha = \varliminf_{n\to+\infty} \alpha \leqslant \varliminf_{n\to+\infty} a_n \leqslant \varlimsup_{n\to+\infty} a_n \leqslant \varlimsup_{n\to+\infty} \beta = \beta.
\]

(2) 设 $\varlimsup\limits_{n\to+\infty} a_n = a_{\text{大}}$，$\varlimsup\limits_{n\to+\infty} b_n = b_{\text{大}}$。由定理~\ref{thm:1.5.3} 知，$\forall \varepsilon > 0$, $\exists N \in \mathbb{N}$，当 $n > N$ 时，有
\[
a_n < a_{\text{大}} + \frac{\varepsilon}{2},\quad b_n < b_{\text{大}} + \frac{\varepsilon}{2},
\]
两式相加得
\[
a_n + b_n < a_{\text{大}} + b_{\text{大}} + \varepsilon.
\]
再由上极限的保不等式性，得
\[
\varlimsup_{n\to+\infty} (a_n + b_n) \leqslant a_{\text{大}} + b_{\text{大}} + \varepsilon.
\]
令 $\varepsilon \to 0^+$，得到
\[
\varlimsup_{n\to+\infty} (a_n + b_n) \leqslant a_{\text{大}} + b_{\text{大}} = \varlimsup_{n\to+\infty} a_n + \varlimsup_{n\to+\infty} b_n.
\]
例如：$a_n = (-1)^n, b_n = (-1)^{n-1}$，则
\[
\varlimsup_{n\to+\infty} (a_n + b_n) = \varlimsup_{n\to+\infty} 0 = 0 < 1 + 1 = \varlimsup_{n\to+\infty} a_n + \varlimsup_{n\to+\infty} b_n.
\]
\end{proof}

\begin{theorem}\label{thm:1.5.6}
\textbf{（上、下极限的等价表达）}
\[
a_{\text{大}} = \varlimsup_{n\to+\infty} x_n = \lim_{n\to+\infty} \sup_{k\geqslant n} \{x_k\} \xlongequal{\text{记作}} \varlimsup_{n\to+\infty} x_n;
\]
\[
a_{\text{小}} = \varliminf_{n\to+\infty} x_n = \lim_{n\to+\infty} \inf_{k\geqslant n} \{x_k\} \xlongequal{\text{记作}} \varliminf_{n\to+\infty} x_n.
\]
因此可用 $\lim\limits_{n\to+\infty} \sup\limits_{k\geqslant n}\{x_k\}$ 与 $\lim\limits_{n\to+\infty} \inf\limits_{k\geqslant n}\{x_k\}$ 分别定义数列 $\{x_n\}$ 的上极限与下极限。
\end{theorem}

\begin{proof}
如果 $\lim\limits_{n\to+\infty} \sup\limits_{k\geqslant n}\{x_k\} = +\infty$，由于 $\sup\limits_{k\geqslant n}\{x_k\}$ 关于 $n$ 单调减，所以 $\sup\limits_{k\geqslant n}\{x_k\} = +\infty$, $\forall n \in \mathbb{N}$。于是，可取 $n_1 \in \mathbb{N}$，s.t. $x_{n_1} > 1$。又可取 $n_2 \in \mathbb{N}$, $n_2 > n_1$，s.t. $x_{n_2} > 2$, $\cdots$。所以，得到数列 $\{x_n\}$ 的子列 $\{x_{n_k}\} \to +\infty\;(k \to +\infty)$。这就证明了 $+\infty$ 为数列 $\{x_n\}$ 的聚点，且为最大聚点。由此得到
\[
\varlimsup_{n\to+\infty} x_n = +\infty = \lim_{n\to+\infty} \sup_{k\geqslant n} \{x_k\}.
\]

如果 $\lim\limits_{n\to+\infty} \sup\limits_{k\geqslant n}\{x_k\} < +\infty$，则 $\lim\limits_{n\to+\infty} \sup\limits_{k\geqslant n}\{x_k\} = -\infty$ 或实数。

设 $a$ 为数列 $\{x_n\}$ 的任一聚点，则必有 $\{x_n\}$ 的子列 $x_{n_i} \to a\;(i \to +\infty)$。$\forall n \in \mathbb{N}$，当 $i \geqslant n$ 时，$n_i \geqslant i \geqslant n$，有
\[
x_{n_i} \leqslant \sup_{k\geqslant n} \{x_k\},
\]
\[
a = \lim_{i\to+\infty} x_{n_i} \leqslant \sup_{k\geqslant n} \{x_k\},
\]
\[
a \leqslant \lim_{n\to+\infty} \sup_{k\geqslant n} \{x_k\},
\]
所以，数列 $\{x_n\}$ 的最大聚点满足
\[
\varlimsup_{n\to+\infty} x_n \leqslant \lim_{n\to+\infty} \sup_{k\geqslant n} \{x_k\}.
\]

另一方面，$\forall y > \varlimsup\limits_{n\to+\infty} x_n$。易见，$[y, +\infty)$ 中最多含数列 $\{x_n\}$ 中的有限项（参阅定理~\ref{thm:1.5.3}(1)）。因此，$\exists N \in \mathbb{N}$，当 $k > N$ 时，有 $x_k < y$。从而，当 $n > N$ 时，有
\[
\sup_{k\geqslant n} \{x_k\} \leqslant y,
\]
由此得到
\[
\lim_{n\to+\infty} \sup_{k\geqslant n} \{x_k\} \leqslant y.
\]
令 $y \to \left(\varlimsup\limits_{n\to+\infty} x_n\right)^+$，推出
\[
\lim_{n\to+\infty} \sup_{k\geqslant n} \{x_k\} \leqslant \varlimsup_{n\to+\infty} x_n.
\]
综合上述，有
\[
\varlimsup_{n\to+\infty} x_n = \lim_{n\to+\infty} \sup_{k\geqslant n} \{x_k\}.
\]
类似证明，或应用上式于 $\{-x_n\}$ 可证得
\[
\varliminf_{n\to+\infty} x_n = \lim_{n\to+\infty} \inf_{k\geqslant n} \{x_k\}.
\]
\end{proof}

\begin{example}\label{ex:1.5.3}
设数列 $\{a_n\}$ 为实数列。$\forall n, m \in \mathbb{N}$，有
\[
a_n \geqslant 0,\quad a_{m+n} \leqslant a_m + a_n.
\]
证明：数列 $\left\{\frac{a_n}{n}\right\}$ 收敛。
\end{example}

\begin{proof}
设 $k$ 是任一固定的自然数，则 $\forall n \in \mathbb{N}, n \geqslant k$，有
\[
n = mk + l,\quad l = 0, 1, \cdots, k-1,\quad m \in \mathbb{N},
\]
所以，
\[
\frac{a_n}{n} = \frac{a_{mk+l}}{n} \leqslant \frac{a_{mk}}{n} + \frac{a_l}{n} \leqslant \frac{m a_k}{n} + \frac{a_l}{n} \leqslant \frac{a_k}{k} + \frac{a_l}{n},
\]
再令 $n \to +\infty$（固定 $k$），可得
\[
\varlimsup_{n\to+\infty} \frac{a_n}{n} \leqslant \lim_{n\to+\infty} \left(\frac{a_k}{k} + \frac{a_l}{n}\right) = \frac{a_k}{k} + 0 = \frac{a_k}{k}.
\]
由此推得
\[
\varlimsup_{n\to+\infty} \frac{a_n}{n} \leqslant \varliminf_{k\to+\infty} \frac{a_k}{k}.
\]
根据定理~\ref{thm:1.5.2} 可知，数列 $\left\{\frac{a_n}{n}\right\}$ 有极限（实数或 $+\infty$）。

又因 $\forall n \in \mathbb{N}$，有
\[
0 \leqslant \frac{a_n}{n} \leqslant \frac{n a_1}{n} = a_1,
\]
所以，
\[
\lim_{n\to+\infty} \frac{a_n}{n} = \varlimsup_{n\to+\infty} \frac{a_n}{n} = \varliminf_{n\to+\infty} \frac{a_n}{n}
\]
为实数，即数列 $\left\{\frac{a_n}{n}\right\}$ 收敛。
\end{proof}

\textbf{广开思路}：请读者想一想，是否能找到其他证明 $\left\{\frac{a_n}{n}\right\}$ 收敛的方法。

\begin{example}\label{ex:1.5.4}
设 $y_n = x_{n-1} + 2x_n$，数列 $\{y_n\}$ 收敛。证明：数列 $\{x_n\}$ 也收敛，且 $\lim\limits_{n\to+\infty} x_n = \frac{1}{3}\lim\limits_{n\to+\infty} y_n$。
\end{example}

\begin{proof}
\textit{证法 1}\quad 设数列 $\{y_n\}$ 收敛于 $y$，则 $\forall \varepsilon > 0$, $\exists N_1 \in \mathbb{N}$，当 $n > N_1$ 时，$|y_n - y| < \frac{\varepsilon}{2}$。于是，当 $n > N > N_1 + \log_2 \frac{\left|x_{N_1} - \frac{y}{3}\right| + 1}{\varepsilon} + 1$ 时，有
\begin{align*}
\left|x_n - \frac{y}{3}\right| &= \left|\left(\frac{y_n}{2} - \frac{x_{n-1}}{2}\right) - \frac{y}{3}\right|\\
&= \left|\frac{1}{2}(y_n - y) - \frac{1}{2}\left(x_{n-1} - \frac{y}{3}\right)\right|\\
&\leqslant \frac{1}{2}|y_n - y| + \frac{1}{2}\left|x_{n-1} - \frac{y}{3}\right|\\
&\leqslant \frac{\varepsilon}{2^2} + \frac{1}{2^2}|y_{n-1} - y| + \frac{1}{2^2}\left|x_{n-2} - \frac{y}{3}\right|\\
&\leqslant \frac{\varepsilon}{2^2} + \frac{\varepsilon}{2^3} + \cdots + \frac{\varepsilon}{2^{n-N_1+1}} + \frac{1}{2^{n-N_1}}\left|x_{N_1} - \frac{y}{3}\right|\\
&< \frac{\varepsilon}{2} + \frac{\varepsilon}{2} = \varepsilon,
\end{align*}
所以，$\{x_n\}$ 收敛于 $\frac{y}{3}$。

\textit{证法 2}\quad 由 $\{y_n\}$ 收敛知，$|y_n| \leqslant M$。因为 $x_n = \frac{y_n}{2} - \frac{x_{n-1}}{2}$，所以
\begin{align*}
|x_n| &= \left|\frac{y_n}{2} - \frac{x_{n-1}}{2}\right| \leqslant \frac{|y_n|}{2} + \frac{|x_{n-1}|}{2}\\
&\leqslant \frac{|y_n|}{2} + \frac{1}{2}\left(\frac{|y_{n-1}|}{2} + \frac{|x_{n-2}|}{2}\right)\\
&\leqslant \cdots \leqslant \left(\frac{|y_n|}{2} + \frac{|y_{n-1}|}{2^2} + \frac{|y_{n-2}|}{2^3} + \cdots\right) + \frac{|x_1|}{2^{n-1}}\\
&\leqslant M\left(\frac{1}{2} + \frac{1}{2^2} + \frac{1}{2^3} + \cdots\right) + |x_1|.
\end{align*}
于是，$\{x_n\}$ 有界，从而 $\varlimsup\limits_{n\to+\infty} x_n$，$\varliminf\limits_{n\to+\infty} x_n$ 均为实数。再根据练习题~1.5 之第~2 题，有
\begin{align*}
\varlimsup_{n\to+\infty} x_n &= \varlimsup_{n\to+\infty} \left(\frac{y_n}{2} - \frac{x_{n-1}}{2}\right) = \lim_{n\to+\infty} \frac{y_n}{2} + \varlimsup_{n\to+\infty} \left(-\frac{x_{n-1}}{2}\right)\\
&= \frac{1}{2}\lim_{n\to+\infty} y_n - \frac{1}{2}\varliminf_{n\to+\infty} x_{n-1}.
\end{align*}
同理，
\begin{align*}
\varliminf_{n\to+\infty} x_n &= \varliminf_{n\to+\infty} \left(\frac{y_n}{2} - \frac{x_{n-1}}{2}\right) = \lim_{n\to+\infty} \frac{y_n}{2} + \varliminf_{n\to+\infty} \left(-\frac{x_{n-1}}{2}\right)\\
&= \frac{1}{2}\lim_{n\to+\infty} y_n - \frac{1}{2}\varlimsup_{n\to+\infty} x_{n-1}.
\end{align*}
整理得到
\[
\begin{cases}
2\varlimsup\limits_{n\to+\infty} x_n = \lim\limits_{n\to+\infty} y_n - \varliminf\limits_{n\to+\infty} x_{n-1},\\[2mm]
2\varliminf\limits_{n\to+\infty} x_n = \lim\limits_{n\to+\infty} y_n - \varlimsup\limits_{n\to+\infty} x_{n-1}.
\end{cases}
\]
两式相减得
\begin{align*}
2\left(\varlimsup_{n\to+\infty} x_n - \varliminf_{n\to+\infty} x_n\right) &= \varlimsup_{n\to+\infty} x_{n-1} - \varliminf_{n\to+\infty} x_{n-1}\\
&= \varlimsup_{n\to+\infty} x_n - \varliminf_{n\to+\infty} x_n,
\end{align*}
即
\[
\varlimsup_{n\to+\infty} x_n - \varliminf_{n\to+\infty} x_n = 0,
\]
\[
\varlimsup_{n\to+\infty} x_n = \varliminf_{n\to+\infty} x_n.
\]
这就证明了数列 $\{x_n\}$ 是收敛的。于是，有
\[
\lim_{n\to+\infty} y_n = \lim_{n\to+\infty} (x_{n-1} + 2x_n) = 3\lim_{n\to+\infty} x_n,
\]
\[
\lim_{n\to+\infty} x_n = \frac{1}{3}\lim_{n\to+\infty} y_n.
\]

\textit{证法 3}\quad 根据练习题~1.5 之第~2 题，有
\begin{align*}
\varliminf_{n\to+\infty} x_n + 2\varlimsup_{n\to+\infty} x_n &= \varliminf_{n\to+\infty} x_{n-1} + \varlimsup_{n\to+\infty} 2x_n \leqslant \varlimsup_{n\to+\infty} (x_{n-1} + 2x_n)\\
&= \varlimsup_{n\to+\infty} y_n = \varliminf_{n\to+\infty} y_n = \varliminf_{n\to+\infty} (x_{n-1} + 2x_n)\\
&\leqslant \varlimsup_{n\to+\infty} x_{n-1} + \varliminf_{n\to+\infty} 2x_n = \varlimsup_{n\to+\infty} x_n + 2\varliminf_{n\to+\infty} x_n,
\end{align*}
所以，
\[
\varliminf_{n\to+\infty} x_n = \varlimsup_{n\to+\infty} x_n.
\]
由证法 1 知，数列 $\{x_n\}$ 有界，故 $\{x_n\}$ 收敛。

\textbf{广开思路}：读者想一想，能否不用上、下极限证明 $\{x_n\}$ 收敛，且 $\lim\limits_{n\to+\infty} x_n = \frac{1}{3}\lim\limits_{n\to+\infty} y_n$。

\textit{证法 4}\quad 易见，
\begin{align*}
& y_1 - 2y_2 + 2^2 y_3 - 2^3 y_4 + \cdots + (-2)^{n-1} y_n\\
=&\; y_1 + \sum_{k=2}^n (-2)^{k-1} y_k = y_1 + \sum_{k=2}^n (-2)^{k-1} (2x_k + x_{k-1})\\
=&\; y_1 - \sum_{k=2}^n (-2)^k x_k + \sum_{k=2}^n (-2)^{k-1} x_{k-1}\\
=&\; y_1 - \sum_{k=2}^n (-2)^k x_k + \sum_{k=1}^{n-1} (-2)^k x_k\\
=&\; y_1 - 2x_1 - (-2)^n x_n,
\end{align*}
所以，
\begin{align*}
\lim_{n\to+\infty} x_n &= \lim_{n\to+\infty} \frac{2x_1 + \sum\limits_{k=2}^n (-2)^{k-1} y_k}{-(-2)^n}\\
&\xlongequal[\text{Stolz 公式}]{\text{定理 1.6.1}} \lim_{n\to+\infty} \frac{(-2)^{n-1} y_n}{-(-2)^n + (-2)^{n-1}}\\
&= \lim_{n\to+\infty} \frac{y_n}{-(-2) + 1} = \lim_{n\to+\infty} \frac{y_n}{3} = \frac{y}{3}.
\end{align*}
\end{proof}

\begin{exblock}{练习题 1.5}
    \item 求 $\varlimsup\limits_{n\to+\infty} a_n$ 与 $\varliminf\limits_{n\to+\infty} a_n$:
    \begin{enumerate}
        \item $a_n = \frac{(-1)^n}{n} + \frac{1 + (-1)^n}{2}$;
        \item $a_n = n^{(-1)^n}$;
        \item $a_n = \left[1 + 2^{(-1)^n n}\right]^{\frac{1}{n}}$;
        \item $a_n = \frac{n^2}{1+n^2}\cos\frac{2n\pi}{3}$;
        \item $a_n = \frac{n^2+1}{n^2}\sin\frac{\pi}{n}$;
        \item $a_n = \sqrt[n]{\left|\cos\frac{n\pi}{3}\right|}$;
        \item $a_n = \begin{cases} 0, & n \text{ 为奇数},\\ \frac{n}{\sqrt[n]{n!}}, & n \text{ 为偶数}. \end{cases}$
    \end{enumerate}

    \item 证明下面各式当两端有意义时成立：
    \begin{enumerate}
        \item 
        \[
        \varliminf_{n\to+\infty} a_n + \varliminf_{n\to+\infty} b_n \leqslant \varliminf_{n\to+\infty} (a_n + b_n) \leqslant \varliminf_{n\to+\infty} a_n + \varlimsup_{n\to+\infty} b_n,
\]
        \[
        \varliminf_{n\to+\infty} a_n + \varlimsup_{n\to+\infty} b_n \leqslant \varlimsup_{n\to+\infty} (a_n + b_n) \leqslant \varlimsup_{n\to+\infty} a_n + \varlimsup_{n\to+\infty} b_n;
\]
        \item 设 $\lim\limits_{n\to+\infty} b_n = b$，则
        \[
        \varliminf_{n\to+\infty} (a_n + b_n) = \varliminf_{n\to+\infty} a_n + b,
\]
        \[
        \varlimsup_{n\to+\infty} (a_n + b_n) = \varlimsup_{n\to+\infty} a_n + b;
\]
        \item 
        \[
        \varliminf_{n\to+\infty} (-a_n) = -\varlimsup_{n\to+\infty} a_n,\quad \varlimsup_{n\to+\infty} (-a_n) = -\varliminf_{n\to+\infty} a_n;
\]
        \item 设 $\{a_n\}$ 与 $\{b_n\}$ 均为非负数列，则
        \[
        \varliminf_{n\to+\infty} a_n \cdot \varliminf_{n\to+\infty} b_n \leqslant \varliminf_{n\to+\infty} a_n b_n \leqslant \varliminf_{n\to+\infty} a_n \cdot \varlimsup_{n\to+\infty} b_n,
\]
        \[
        \varliminf_{n\to+\infty} a_n \cdot \varlimsup_{n\to+\infty} b_n \leqslant \varlimsup_{n\to+\infty} a_n b_n \leqslant \varlimsup_{n\to+\infty} a_n \cdot \varlimsup_{n\to+\infty} b_n;
\]
        \item 设 $\{b_n\}$ 非负，且 $\lim\limits_{n\to+\infty} b_n = b$，则
        \[
        \varliminf_{n\to+\infty} a_n b_n = b \varliminf_{n\to+\infty} a_n,\quad \varlimsup_{n\to+\infty} a_n b_n = b \varlimsup_{n\to+\infty} a_n;
\]
        \item 设 $a_n > 0\;(n \in \mathbb{N})$，$\varliminf\limits_{n\to+\infty} a_n > 0$，则
        \[
        \varliminf_{n\to+\infty} \frac{1}{a_n} = \frac{1}{\varlimsup\limits_{n\to+\infty} a_n}.
\]
    \end{enumerate}

    \item 设 $a_n > 0\;(n \in \mathbb{N})$，且 $\varliminf\limits_{n\to+\infty} a_n \cdot \varlimsup\limits_{n\to+\infty} \frac{1}{a_n} = 1$，证明：数列 $\{a_n\}$ 收敛。

    \item 设数列 $\{a_n\}$，$a_n \leqslant 1$，$n = 1, 2, \cdots$，且满足：
    \[
    a_m + a_n - 1 < a_{m+n} < a_m + a_n + 1.
\]
    证明：
    \begin{enumerate}
        \item $\lim\limits_{n\to+\infty} \frac{a_n}{n} = \omega$，其中 $\omega$ 为有限数；
        \item $n\omega - 1 \leqslant a_n \leqslant n\omega + 1$。
    \end{enumerate}
    （提示：将例~\ref{ex:1.5.3} 的结果用到数列 $\{a_n + 1\}$ 与 $\{1 - a_n\}$ 上。）

    \item 设 $a_n \geqslant 0$, $n \in \mathbb{N}$。证明：
    \[
    \varlimsup_{n\to+\infty} \sqrt[n]{a_n} \leqslant 1 \iff \text{对任何 } l > 1\text{，有 } \lim_{n\to+\infty} \frac{a_n}{l^n} = 0.
\]
    如果删去“任何”两字，结论如何？
\end{exblock}

\begin{exblock}{思考题 1.5}
    \item 设数列 $\{x_n\}$ 有界，且 $\lim\limits_{n\to+\infty} (x_{n+1} - x_n) = 0$，令
    \[
    l = \varliminf_{n\to+\infty} x_n,\quad L = \varlimsup_{n\to+\infty} x_n.
\]
    证明：$\{a \in \mathbb{R} \mid \text{有子列 } x_{n_k} \to a\;(k \to \infty)\} = [l, L]$。如果删去条件 $\lim\limits_{n\to+\infty} (x_{n+1} - x_n) = 0$，结论如何？

    \item 设 $0 \leqslant a_{n+m} \leqslant a_n \cdot a_m\;(n, m = 1, 2, \cdots)$。证明：$\varlimsup\limits_{n\to+\infty} \sqrt[n]{a_n} = \varliminf\limits_{n\to+\infty} \sqrt[n]{a_n}$ 且 $\sqrt[n]{a_n}$ 收敛。
\end{exblock}
